\section{The block transfer}\label{sec:transfer}

For a partition $\nu$, let
\[
 \operatorname{supp}(\nu)=\{a\geq1:m_a(\nu)>0\},
\]
where $m_a(\nu)$ is the multiplicity of the part $a$. The transfer below
changes multiplicities, while the block decompositions are determined by this
finite support set.

\begin{lemma}[Greedy interval lemma]\label{lem:greedy-interval}
Let $D$ be a finite set of integers and let $k\geq1$. Partition $D$ by the
following two greedy procedures:
\begin{enumerate}
\item start with the least unused element $b$, take all unused elements in
$[b,b+k]$, and continue upward;
\item start with the greatest unused element $a$, take all unused elements in
$[a-k,a]$, and continue downward.
\end{enumerate}
The two procedures produce the same number of blocks.
\end{lemma}

\begin{proof}
In the first procedure, record the least element selected at each step. These
recorded elements differ pairwise by more than $k$, and every subset of $D$
whose distinct elements differ pairwise by more than $k$ contains at most one
element in each greedy block. Thus the number of blocks in the first
procedure is the maximum cardinality of a $k$-separated subset of $D$.
The greatest-first procedure has the same property, using its greatest
elements as the recorded separated subset. Hence the two block counts agree.
\end{proof}

Let $\lambda\in\Rk$ be nonempty. Define $\tau(\lambda)$ by acting
simultaneously on its $k$-blocks. From every block with largest part $a>k$,
replace one occurrence of $a$ by $a-k$. From the final block with largest
part $a\leq k$, if it exists, delete one occurrence of $a$. The resulting
partition is again $k$-regular. The decrease in size is the sum of the block
weights:
\begin{equation}\label{eq:size-transfer}
 |\tau(\lambda)|=|\lambda|-\bk(\lambda).
\end{equation}

The reverse operation is explicit. For an integer $w>0$, let $r$ be its least
nonnegative residue modulo $k$.

\begin{lemma}[Explicit reverse transfer]\label{lem:block-transfer}
Let $\mu\in\Rk$, put $h=\bk(\mu)$, and suppose that
\begin{equation}\label{eq:admissible-weight}
 h\leq w\leq h+k-1.
\end{equation}
Construct $\sigma_w(\mu)$ by the following finite algorithm.
\begin{enumerate}
\item If $r>0$, adjoin one part $r$ to $\mu$.
\item Leave all parts of size at most $r$ fixed. On the distinct values in
$\operatorname{supp}(\mu)$ larger than $r$, work from the smallest upward. If
$b$ is the smallest unprocessed value, make a reverse block from all
unprocessed values in $[b,b+k]$, and continue with values larger than $b+k$.
\item In every reverse block, replace one occurrence of its smallest part $b$
by $b+k$.
\end{enumerate}
Then $\sigma_w(\mu)\in\Rk$ and
\begin{equation}\label{eq:transfer-inverse}
 \tau\bigl(\sigma_w(\mu)\bigr)=\mu,
 \qquad
 \bk\bigl(\sigma_w(\mu)\bigr)=w.
\end{equation}
\end{lemma}

\begin{proof}
The algorithm is deterministic and terminates because the support is finite.
Every new value has the same nonzero residue modulo $k$ as the value it
replaces, so $\sigma_w(\mu)$ is $k$-regular.

Write
\begin{equation}\label{eq:decompose-h}
 h=ak+t,\qquad 0\leq t<k.
\end{equation}
There are $a$ blocks of $\mu$ whose largest parts exceed $k$; if $t>0$, the
remaining block has largest part $t$. Let
\[
 D_r=\{d\in\operatorname{supp}(\mu):d>r\}.
\]
Removing the values at most $r$ removes the terminal block when $r\geq t$,
and leaves a nonempty part of it when $r<t$. The blocks with largest parts
greater than $k$ retain their largest parts, so the greatest-first
decomposition of $D_r$ has $a$ blocks when $r\geq t$, and $a+1$ blocks when
$r<t$. By \Cref{lem:greedy-interval}, the reverse-block algorithm in step 2
has the same number, say $q$, of blocks. Consequently $q=a$ in the first
case and $q=a+1$ in the second. The interval condition
\eqref{eq:admissible-weight} gives $w=ak+r$ in the first case and
$w=(a+1)k+r$ in the second. Therefore $qk+r=w$.

Let the reverse blocks have smallest values $b_1<\cdots<b_q$. Consecutive
smallest values satisfy $b_{j+1}>b_j+k$. After step 3, the values originating
in the $j$th reverse block lie in $[b_j,b_j+k]$ and include $b_j+k$.
Moreover, $b_j+k<b_{j+1}$, so the top-down decomposition cannot merge two
adjacent reverse blocks. It therefore recovers these intervals, in reverse
order, as blocks with largest values $b_q+k,\ldots,b_1+k$. If $r>0$, all
fixed values are at most $r$ and the adjoined value $r$ is present. When
$q>0$, they lie below $b_1$ and form the final block; when $q=0$, they form
the only block. In both cases its largest value is $r$.

Thus the block weights of $\sigma_w(\mu)$ are $q$ copies of $k$ and, when
$r>0$, one weight $r$; their sum is $qk+r=w$. Applying $\tau$ lowers one
$b_j+k$ to $b_j$ in every nonterminal block and removes the adjoined part
$r$ from the terminal block. All multiplicities are thereby restored, so
$\tau(\sigma_w(\mu))=\mu$.
\end{proof}

\begin{lemma}[Forward transfer is inverted by the algorithm]
\label{lem:converse-transfer}
Let $\lambda\in\Rk$ be nonempty, put $w=\bk(\lambda)$, and set
$\mu=\tau(\lambda)$. Then
\begin{equation}\label{eq:block-drop}
 \bk(\mu)\leq w\leq\bk(\mu)+k-1,
\end{equation}
and $\lambda=\sigma_w(\mu)$.
\end{lemma}

\begin{proof}
Write $w=ak+r$ with $0\leq r<k$. The $a$ nonterminal blocks of $\lambda$ have
largest parts $A_1>\cdots>A_a>k$. If $r>0$, the final block has largest
part $r$; if $r=0$, there is no final block of weight less than $k$.

For $1\leq j\leq a$, put $c_j=A_j-k$. The image under $\tau$ of the $j$th
nonterminal block has support contained in $[c_j,c_j+k]$ and contains
$c_j$, because one occurrence of $A_j$ was moved to $A_j-k$. For
$1\leq j<a$, the next lower nonterminal block satisfies
\[
 c_{j+1}+k=A_{j+1}<A_j-k=c_j,
\]
by the defining greedy separation of the original blocks. Thus the support
intervals $[c_j,c_j+k]$ are pairwise separated. If a terminal block exists,
its largest part is $r<c_a$; after the deletion, every surviving terminal
value is at most $r$.

Consequently, the distinct values of $\mu$ larger than $r$ are partitioned by
the reverse greedy procedure into exactly the $a$ transformed nonterminal
blocks. It starts at $c_a$, recovers the support in $[c_a,c_a+k]$, and then
proceeds successively through the separated intervals up to $[c_1,c_1+k]$.
The values at most $r$ are precisely the unchanged terminal values. Step 3 of
the construction of $\sigma_w(\mu)$ therefore reverses every move made by
$\tau$, including the possible terminal deletion. Hence $\lambda=\sigma_w(\mu)$.

It remains to compare block indices. Set
\[
 D_r=\{d\in\operatorname{supp}(\mu):d>r\}.
\]
If $D_r$ is empty, then $a=0$ and all parts of $\mu$ are at most $r$.
Hence $\bk(\mu)=t$ for some $0\leq t\leq r$, where $t=0$ means that $\mu$
is empty.

Suppose that $D_r$ is nonempty. Its bottom-up decomposition has the $a$
transformed nonterminal blocks identified above, so \Cref{lem:greedy-interval}
shows that its top-down decomposition also has exactly $a$ blocks. Since $\mu$
is $k$-regular, the top of its lowest block cannot equal $k$. If that top is
greater than $k$, then all $a$ blocks have weight $k$. Values at most $r$ either
join the lowest block or form a terminal block of top $t\leq r$; in either event
\[
 \bk(\mu)=ak+t\qquad(0\leq t\leq r),
\]
with $t=0$ when there is no terminal block. If the lowest block has top less
than $k$, it is itself the terminal block and contains every positive value
of $\mu$ at most $r$. Its top $t$ satisfies $r<t<k$, while the other $a-1$
blocks have weight $k$. Thus
\[
 \bk(\mu)=(a-1)k+t\qquad(r<t<k).
\]
The first formula also covers $D_r=\varnothing$. In the two cases,
respectively,
\[
 w-\bk(\mu)=r-t,
 \qquad
 w-\bk(\mu)=k+r-t.
\]
Both differences lie in $\{0,1,\ldots,k-1\}$, which proves
\eqref{eq:block-drop}.
\end{proof}
